\documentclass[aspectratio=169,11pt]{beamer}

% ── Colores ──────────────────────────────────────────────
\definecolor{azulnoche}{RGB}{10,15,45}
\definecolor{azulmedio}{RGB}{20,40,90}
\definecolor{azulclaro}{RGB}{40,80,160}
\definecolor{blanco}{RGB}{255,255,255}
\definecolor{acento}{RGB}{100,160,255}

\usetheme{default}
\usecolortheme{default}

\setbeamercolor{background canvas}{bg=azulnoche}
\setbeamercolor{normal text}{fg=blanco}
\setbeamercolor{frametitle}{fg=blanco,bg=azulmedio}
\setbeamercolor{title}{fg=blanco}
\setbeamercolor{subtitle}{fg=acento}
\setbeamercolor{author}{fg=acento}
\setbeamercolor{institute}{fg=acento!70!white}
\setbeamercolor{date}{fg=acento!60!white}
\setbeamercolor{item}{fg=acento}
\setbeamercolor{item projected}{fg=blanco,bg=azulclaro}
\setbeamercolor{block title}{fg=blanco,bg=azulclaro}
\setbeamercolor{block body}{fg=blanco,bg=azulmedio!60}
\setbeamercolor{block title example}{fg=blanco,bg=azulclaro!80}
\setbeamercolor{block body example}{fg=blanco,bg=azulmedio!50}
\setbeamercolor{structure}{fg=acento}
\setbeamercolor{section in toc}{fg=acento}
\setbeamercolor{subsection in toc}{fg=acento!70!white}

\setbeamertemplate{navigation symbols}{}
\setbeamertemplate{footline}{\hfill\insertframenumber/\inserttotalframenumber\hspace*{8pt}\vspace{4pt}}
\setbeamertemplate{itemize items}[circle]

\usepackage[utf8]{inputenc}
\usepackage[T1]{fontenc}
\usepackage{amsmath,amssymb}
\usepackage{graphicx}
\usepackage{booktabs}
\usepackage{tikz}
\usetikzlibrary{shapes,arrows.meta,positioning,calc}

% ── Documento ────────────────────────────────────────────
\title{Predictor Mundial 2026}
\subtitle{Modelo de predicci\'on basado en aprendizaje autom\'atico}
\author{Fernando Reyna Rodr\'{\i}guez}
\institute{Universidad Toribio Rodr\'{\i}guez de Mendoza\\Procesamiento de Lenguaje Natural}
\date{\today}

\begin{document}

% ═══════════════════════════════════════════════════════════
% PORTADA
% ═══════════════════════════════════════════════════════════
\begin{frame}[plain]
\begin{tikzpicture}[remember picture,overlay]
  \fill[azulmedio] (current page.north west) rectangle ([yshift=-4cm]current page.north east);
  \fill[azulclaro!30] ([yshift=-4cm]current page.north west) rectangle ([yshift=-4.15cm]current page.north east);
\end{tikzpicture}

\vspace*{1.5cm}
\begin{center}
  {\Huge\bfseries Predictor Mundial 2026}\\[12pt]
  {\Large\color{acento} Modelo de predicci\'on basado en aprendizaje autom\'atico}\\[30pt]
  {\large\color{blanco} Fernando Reyna Rodr\'{\i}guez}\\[8pt]
  {\color{acento} Universidad Toribio Rodr\'{\i}guez de Mendoza}\\[4pt]
  {\color{acento!70!white} Procesamiento de Lenguaje Natural}\\[20pt]
  {\small\color{acento!60!white} \today}
\end{center}
\end{frame}

% ═══════════════════════════════════════════════════════════
% OBJETIVO
% ═══════════════════════════════════════════════════════════
\begin{frame}{Objetivo del Proyecto}
\begin{itemize}
  \item Construir un pipeline de predicci\'on completo para el Mundial FIFA 2026.
  \item Utilizar \textbf{aprendizaje autom\'atico} para estimar probabilidades de victoria.
  \item Simular el torneo completo (fase de grupos y knockout) con \textbf{48 selecciones}.
  \item Generar un dashboard interactivo para explorar resultados.
\end{itemize}

\vspace{12pt}
\begin{block}{Idea central}
Combinar predicciones de un modelo ML con el sistema de \textbf{Elo} para obtener probabilidades m\'as robustas en cada partido.
\end{block}
\end{frame}

% ═══════════════════════════════════════════════════════════
% MODELO
% ═══════════════════════════════════════════════════════════
\begin{frame}{Modelo de Predicci\'on}
\begin{columns}[T]
\begin{column}{0.52\textwidth}
\textbf{\color{acento} Ensemble: Random Forest + XGBoost}

\vspace{6pt}
\begin{itemize}
  \item \textbf{Random Forest}: 600 \'arboles, prof. m\'ax. 14
  \item \textbf{XGBoost}: 400 \'arboles, prof. 6, ratio 0.05
  \item Promedio de probabilidades de ambos modelos
  \item Divisi\'on: 80\% train / 20\% test
\end{itemize}
\end{column}

\begin{column}{0.45\textwidth}
\begin{block}{M\'etricas}
\begin{tabular}{lr}
\toprule
\textbf{M\'etrica} & \textbf{Valor} \\
\midrule
Accuracy    & 67.9\%  \\
Precision   & 68.2\%  \\
Recall      & 64.7\%  \\
F1-Score    & 0.664   \\
ROC-AUC     & 0.752   \\
\bottomrule
\end{tabular}
\end{block}
\end{column}
\end{columns}
\end{frame}

% ═══════════════════════════════════════════════════════════
% FEATURES
% ═══════════════════════════════════════════════════════════
\begin{frame}{Variables del Modelo (Features)}
\begin{columns}[T]
\begin{column}{0.5\textwidth}
\textbf{\color{acento} 10 variables de entrada:}
\vspace{4pt}
\begin{enumerate}
  \item \texttt{elo\_diff} --- Diferencia de Elo
  \item \texttt{form\_diff} --- Diferencia de forma reciente
  \item \texttt{rank\_diff} --- Diferencia de ranking
  \item \texttt{home\_advantage} --- Ventaja de local
  \item \texttt{avg\_rating} --- Calidad promedio del partido
\end{enumerate}
\end{column}

\begin{column}{0.5\textwidth}
\vspace{22pt}
\begin{enumerate}
  \setcounter{enumi}{5}
  \item \texttt{rating\_ratio} --- Fuerza relativa local
  \item \texttt{elo\__squared} --- Elo al cuadrado
  \item \texttt{abs\_elo\_diff} --- Brecha absoluta Elo
  \item \texttt{form\_abs\_diff} --- Brecha absoluta forma
  \item \texttt{elo\_power} --- Elo escalado (tanh)
\end{enumerate}
\end{column}
\end{columns}

\vspace{10pt}
\begin{block}{Variable objetivo}
\texttt{target = 1} si el equipo local gana, \texttt{0} en otro caso.\\
\small{\color{acento!70!white} Datos desde 2018 $\cdot$ Ratings Elo de 244 equipos $\cdot$ Rankings FIFA de 211 equipos}
\end{block}
\end{frame}

% ═══════════════════════════════════════════════════════════
% SIMULACION DEL TORNEO
% ═══════════════════════════════════════════════════════════
\begin{frame}{Simulaci\'on del Torneo}
\begin{columns}[T]
\begin{column}{0.55\textwidth}
\textbf{\color{acento} Formato Mundial 2026}

\vspace{4pt}
\begin{itemize}
  \item \textbf{48 selecciones} en 12 grupos (A--L)
  \item Sistema serpenteo por clasificaci\'on Elo
  \item 72 partidos en fase de grupos
  \item Eliminaci\'on directa: 32avos $\to$ Final
  \item \textbf{103 partidos por simulaci\'on}
  \item \textbf{20 simulaciones} con semillas distintas
\end{itemize}
\end{column}

\begin{column}{0.42\textwidth}
\begin{block}{F\'ormula de probabilidad}
\vspace{-2pt}
\[ P_{final} = \alpha \cdot P_{ML} + (1-\alpha) \cdot P_{Elo} \]

\vspace{4pt}
{\small\color{acento!70!white}
$\alpha = 0.20$ (knockout)\\
$\alpha = 0.25$ (grupos)\\[4pt]
$P_{Elo} = \dfrac{1}{1 + 10^{-\Delta Elo/400}}$
}
\end{block}
\end{column}
\end{columns}
\end{frame}

% ═══════════════════════════════════════════════════════════
% ESTRUCTURA DEL CODIGO
% ═══════════════════════════════════════════════════════════
\begin{frame}{Estructura del Proyecto}
\begin{center}
\begin{tikzpicture}[
  box/.style={rectangle, rounded corners=3pt, draw=acento, fill=azulmedio,
    text=blanco, minimum height=0.7cm, minimum width=2.2cm, font=\small\bfseries},
  arr/.style={-{Stealth[length=5pt]}, thick, acento}
]
\node[box] (data) at (0,0) {data/};
\node[box] (src) at (3,0) {src/};
\node[box] (models) at (6,0) {models/};
\node[box] (outputs) at (9,0) {outputs/};

\node[box, fill=azulclaro] (app) at (4.5,-1.4) {app.py (Streamlit)};

\draw[arr] (data) -- node[above,font=\tiny,fg=acento]{carga} (src);
\draw[arr] (src) -- node[above,font=\tiny,fg=acento]{entrena} (models);
\draw[arr] (models) -- node[above,font=\tiny,fg=acento]{predice} (outputs);
\draw[arr] (src) -- (app);
\draw[arr] (outputs) -- (app);
\end{tikzpicture}
\end{center}

\vspace{8pt}
\begin{itemize}
  \item \texttt{src/data\_loader.py} --- Carga de datos hist\'oricos, Elo y FIFA
  \item \texttt{src/features.py} --- Construcci\'on de las 10 variables
  \item \texttt{src/models.py} --- Entrenamiento RF + XGBoost
  \item \texttt{src/simulate\_bracket.py} --- Simulaci\'on del torneo completo
  \item \texttt{predict\_world\_cup.py} --- Script principal (orchestrator)
  \item \texttt{app.py} --- Dashboard interactivo con Streamlit
\end{itemize}
\end{frame}

% ═══════════════════════════════════════════════════════════
% RESULTADOS
% ═══════════════════════════════════════════════════════════
\begin{frame}{Resultados Predichos}
\begin{columns}[T]
\begin{column}{0.5\textwidth}
\begin{block}{Campe\'on}
\vspace{4pt}
\begin{center}
{\Large\bfseries \color{blanco} \España}\\[4pt]
{\small 13 de 20 simulaciones (65\%)}
\end{center}
\vspace{-4pt}
\end{block}

\vspace{6pt}
\begin{block}{Subcampe\'on}
\begin{center}
{\Large\bfseries \color{blanco} B\'elgica}\\[2pt]
{\small Final: Espa\~na 2 -- 0 B\'elgica}
\end{center}
\vspace{-4pt}
\end{block}
\end{column}

\begin{column}{0.48\textwidth}
\begin{block}{Frecuencia de campe\'on}
\vspace{4pt}
\begin{tabular}{lcc}
\toprule
\textbf{Selecci\'on} & \textbf{Sim.} & \textbf{\%} \\
\midrule
\España     & 13 & 65\% \\
\Argentina  & 6  & 30\% \\
Francia     & 1  & 5\%  \\
\bottomrule
\end{tabular}
\end{block}

\vspace{6pt}
{\small\color{acento}
\textbf{Semifinalistas:} Espa\~na, Noruega, Ecuador, B\'elgica}
\end{column}
\end{columns}
\end{frame}

% ═══════════════════════════════════════════════════════════
% CONCLUSIONES
% ═══════════════════════════════════════════════════════════
\begin{frame}{Conclusiones}
\begin{itemize}
  \item El modelo \textbf{ensemble} (RF + XGBoost) combinado con \textbf{Elo} logra un accuracy del 67.9\% y ROC-AUC de 0.752.
  \item La combinaci\'on ponderada ML--Elo permite integrar el conocimiento estad\'{\i}stico cl\'asico con el aprendizaje autom\'atico.
  \item El pipeline es \textbf{reproducible, modular y extensible} con datos reales.
  \item El dashboard Streamlit permite explorar cada grupo, fase y partido de forma interactiva.
\end{itemize}

\vspace{10pt}
\begin{block}{Mejoras futuras}
\begin{itemize}
  \item Incorporar datos reales de partidos y lesiones
  \item A\~nadir m\'as features (posesi\'on, tiros, etc.)
  \item Optimizar hiperpar\'ametros con Bayesian Search
\end{itemize}
\end{block}
\end{frame}

% ═══════════════════════════════════════════════════════════
% CIERRE
% ═══════════════════════════════════════════════════════════
\begin{frame}[plain]
\begin{center}
\vspace*{3cm}
{\Huge\bfseries \color{blanco} \'Gracias}\\[12pt]
{\Large\color{acento} Fernando Reyna Rodr\'{\i}guez}\\[6pt]
{\color{acento!70!white} fernando.reyna@ucontrm.edu.pe}
\end{center}
\end{frame}

\end{document}
